\section{总结与延伸}

\subsection{讲者的核心总结}

Tatsu Hashimoto 在课程结尾强调了以下要点：

\begin{enumerate}
    \item \textbf{没有银弹}：单一并行策略无法解决所有问题，需要组合数据并行、张量并行和流水线并行
    \item \textbf{硬件拓扑决定策略}：不同层级的带宽差异直接决定了应在哪个层级使用哪种并行
    \item \textbf{简洁的经验法则}：先用张量并行填满节点、再跨节点扩展、最后数据并行填充剩余 GPU
    \item \textbf{批次大小是资源}：有限的批次大小需要在数据并行度和流水线微批次数之间权衡
\end{enumerate}

\subsection{全课知识图谱}

\begin{center}
\begin{tikzpicture}[node distance=0.8cm, every node/.style={draw, rounded corners, minimum width=3cm, minimum height=0.7cm, font=\small}]
\node (hw) {网络层次与集合通信};
\node (dp) [below=of hw] {数据并行 (DDP/ZeRO/FSDP)};
\node (mp) [below=of dp] {模型并行 (PP + TP)};
\node (sp) [below=of mp] {序列并行 + 激活优化};
\node (combine) [below=of sp] {3D/4D 并行组合};
\draw[->, thick] (hw) -- (dp);
\draw[->, thick] (dp) -- (mp);
\draw[->, thick] (mp) -- (sp);
\draw[->, thick] (sp) -- (combine);
\end{tikzpicture}
\end{center}

\subsection{关键公式汇总}

\begin{itemize}
    \item \textbf{All-Reduce 恒等式}：$\text{All-Reduce} \equiv \text{Reduce-Scatter} + \text{All-Gather}$
    \item \textbf{每参数训练内存}：$16$ 字节（BF16 混合精度 + AdamW）
    \item \textbf{流水线气泡比例}：$(P-1)/M$（$P$ = 阶段数，$M$ = 微批次数）
    \item \textbf{每层激活内存}：$SBH \cdot (34 + 5AS/H)$
    \item \textbf{优化后激活内存}：$34 \cdot SBH / T$（张量+序列并行+重计算）
    \item \textbf{FSDP 通信量}：$3|\theta|$（vs 朴素 DP 的 $2|\theta|$）
\end{itemize}

\subsection{拓展阅读}

\begin{itemize}
    \item Rajbhandari et al., ZeRO: Memory Optimizations Toward Training Trillion Parameter Models: \url{https://arxiv.org/abs/1910.02054}
    \item Narayanan et al., Efficient Large-Scale Language Model Training on GPU Clusters Using Megatron-LM: \url{https://arxiv.org/abs/2104.04473}
    \item Llama 3 Technical Report: \url{https://arxiv.org/abs/2407.21783}
    \item DeepSeek-V3 Technical Report: \url{https://arxiv.org/abs/2412.19437}
    \item Google, Introduction to TPU Training (Parallelism Chapter): \url{https://cloud.google.com/tpu/docs/introduction-to-tpu-training}
    \item Zero Bubble Pipeline Parallelism: \url{https://arxiv.org/abs/2401.10241}
    \item OpenAI, An Empirical Model of Large-Batch Training (Critical Batch Size): \url{https://arxiv.org/abs/1812.06162}
\end{itemize}

\end{document}
